\subsection{APPLICATION TO LINEAR EQUATIONS}

Consider a system of \(n\) scalar linear equations in \(n\) unknowns over a (commutative) ring A (II, § 2, no. 8):

\[
\sum_ {j = 1} ^ {n} \lambda_ {i j} \xi_ {j} = \eta_ {i} \quad (1 \leqslant i \leqslant n).\tag{34}
\]

Let \(L\) be the square matrix \((\lambda_{ij})\) of order \(n\) ; by identifying as usual the matrix with one column consisting of the \(\xi_i\) (resp. the \(\eta_i\) ) with an element \(x = (\xi_i)\) of \(A^n\) (resp. the element \(y = (\eta_i)\) of \(A^n\) ), system (34) may also be written (II, § 10, no. 3, Proposition 2)

\[
L. x = y.\tag{35}
\]

Let \(u\) be the endomorphism \(x \mapsto L.x\) of \(A^n\) , with \(L\) as matrix with respect to the canonical basis; to say that equation (35) has (at least) one solution for all \(y \in A^n\) means that \(u\) is surjective; Theorem 1 of no. 2 then implies the following proposition:

PROPOSITION 13. For a system of n linear equations in n unknowns over a commutative ring to admit at least one solution regardless of the right hand sides, it is necessary and sufficient that the determinant of the matrix of the system be invertible; in this case the system admits a single solution.

If det L is not a divisor of zero in A, equation (34) is equivalent to the equation

\[
(\det L) L. x = (\det L) y.
\]

If \(M\) is the cofactor matrix of \(L\) , we derive from (34) and formula (26) of no. 6 the relation

\[
(\det L) x = ^ {t} M. y\tag{36}
\]

which may also be written

\[
(\det L) \xi_ {i} = \sum_ {j = 1} ^ {n} (- 1) ^ {i + j} (\det L ^ {i j}) \eta_ {j} = \det L _ {i} \quad (1 \leqslant i \leqslant n)\tag{37}
\]

where \(L_{i}\) denotes the matrix obtained by replacing the column of L of index i by y. Formulae (37) are called Cramer's formulae for system (34); every solution of (34) is also a solution of (37). Conversely, we derive from (36), taking account of formula (28) of no. 6,

\[
(\det L) (L. x - y) = 0\tag{38}
\]

and hence, if det L is not a divisor of zero in A, systems (34) and (37) are equivalent; if det L is invertible, the unique solution of (34) is given by

\[
\xi_ {i} = (\det L) ^ {- 1} (\det L _ {i}) \quad (1 \leqslant i \leqslant n).\tag{39}
\]

A system (34) such that det L is invertible is also called a Cramer system.

In particular let \(y = 0\) ; it then follows from Proposition 3 of no. 2 that:

PROPOSITION 14. For a homogeneous linear system of n equations in n unknowns over a commutative ring to admit a non-zero solution, it is necessary and sufficient that the determinant of its matrix be a divisor of zero.

